\subsection{Qué hacía el TP x86}

Este ejercicio era el más denso del TP3. Incluía:

\begin{itemize}
  \item \texttt{mmu\_init}: contadores \texttt{next\_free\_kernel\_page} (desde \texttt{0x100000}) y
  \texttt{next\_free\_task\_page} (desde \texttt{0x400000}).
  \item \texttt{mmu\_initTaskDir}: construir el PD/PT de la tarea, identity-mapear kernel + área libre kernel, copiar el
  código original de la tarea (desde \texttt{MMU\_BASE\_ADDR\_TAREAS}) a su celda física en el tablero, mapear
  \texttt{TASK\_CODE} (\texttt{0x08000000}) a esas páginas con permisos U, y mapear \texttt{TASK\_COMPARTIDA} a
  \texttt{g\_shared[jugador]}.
  \item \texttt{mmu\_mapPage} / \texttt{mmu\_unmapPage}: helpers para mapear/desmapear una página virtual, invalidando
  el TLB con \texttt{tlbflush}.
\end{itemize}

\subsection{Equivalente en RISC-V}

La portación conserva el esquema ``el código del minero vive físicamente en su celda del tablero''. Las funciones
equivalentes viven en \texttt{riscv/mmu.cpp}:

\begin{itemize}
  \item \texttt{mmu\_task\_init(slot, player, x, y)}: equivalente de \texttt{mmu\_initTaskDir}.
  \item \texttt{mmu\_task\_move(slot, old\_x, old\_y, new\_x, new\_y)}: equivalente conjunto de ``copiar 8~KiB +
  remapear \texttt{TASK\_CODE}''.
  \item \texttt{mmu\_switch\_to\_task(slot)} / \texttt{mmu\_switch\_to\_idle()}: análogo a ``cargar CR3'' (escribir
  \texttt{satp}).
\end{itemize}

No se exponen \texttt{mmu\_mapPage}/\texttt{mmu\_unmapPage} genéricas, porque en este port el universo de mapeos es más
restringido y conocido (kernel gigapage, UART megapage, dos páginas de código y dos páginas compartidas por tarea). El
rol de esas funciones lo cumplen \texttt{map\_user\_code\_pages} y \texttt{map\_user\_shared}, que actualizan
directamente \texttt{t.l0\_user[]}.

\subsection{Allocator del kernel}

El análogo del ``contador de páginas libres'' del TP x86 es un \emph{bump allocator} a partir del final del
\texttt{.bss}:

\begin{codesnippet}
\begin{verbatim}
 static uint64_t g_alloc_page = 0;

 static void* alloc_page_table() {
     if (g_alloc_page == 0) {
         g_alloc_page = align_up_u64(
             reinterpret_cast<uint64_t>(&__bss_end), PAGE_SIZE);
     }
     void* page = reinterpret_cast<void*>(g_alloc_page);
     g_alloc_page += PAGE_SIZE;
     memset(page, 0, PAGE_SIZE);
     return page;
 }
\end{verbatim}
\end{codesnippet}

\subsection{``Tablero como memoria''}

El ``área mapa'' del TP3 (\texttt{0x400000}..\texttt{0x1C5FFFF}, 78$\times$40 celdas de 8~KiB) se reemplaza por un
\textbf{arreglo estático alineado a página}:

\begin{codesnippet}
\begin{verbatim}
 alignas(PAGE_SIZE) static uint8_t
     g_board[TABLERO_N_COLS * TABLERO_N_FILS * kCellBytes];

 static inline uint64_t board_cell_pa(uint32_t x, uint32_t y) {
     const uint64_t ix = y * TABLERO_N_COLS + x;
     return reinterpret_cast<uint64_t>(&g_board[ix * kCellBytes]);
 }
\end{verbatim}
\end{codesnippet}

Esto evita tener que fijar direcciones físicas ``mágicas'' y deja que el linker resuelva el layout. Igualmente se cumple
la propiedad del TP3: cada celda del tablero es un bloque contiguo de 8~KiB y \texttt{board\_cell\_pa(x,y)} es su
dirección física.

Las áreas compartidas entre mineros del mismo jugador se resuelven también con arreglos estáticos:

\begin{codesnippet}
\begin{verbatim}
 alignas(PAGE_SIZE) static uint8_t g_shared[2][2 * PAGE_SIZE];
\end{verbatim}
\end{codesnippet}

\subsection{Construcción de la page-table de tarea}

\texttt{mmu\_task\_init} construye la page-table completa de una tarea:

\begin{codesnippet}
\begin{verbatim}
 static inline void task_build_pagetable(TaskMmu& t,
         uint16_t asid, uint32_t player_id, uint64_t cell_pa) {
     t.root = reinterpret_cast<uint64_t*>(alloc_page_table());
     map_kernel_gigapage(t.root);
     ensure_l1_0(t);
     map_uart(t);
     map_user_code_pages(t, cell_pa);       // TASK_CODE (2 pags)
     map_user_shared(t, player_id);         // TASK_COMPARTIDA (2 pags)
     t.asid = asid;
     t.satp = make_satp(
         reinterpret_cast<uint64_t>(t.root), asid);
     t.initialized = 1;
 }
\end{verbatim}
\end{codesnippet}

El mapeo de \texttt{TASK\_CODE} es el punto crítico de la mecánica del juego: a partir de la VA \texttt{0x08000000} se
mapean dos páginas físicas contiguas a partir de \texttt{cell\_pa}, con flags \texttt{V|R|W|X|U|A|D}. Es decir, \emph{el
minero puede leer, escribir y ejecutar en el mismo bloque de memoria} --- exactamente la condición que hace posible que
el juego copie el código de la tarea a otra celda y la tarea siga ejecutando desde \texttt{0x08000000} con código
distinto.

\subsection{\texttt{mmu\_task\_move}: movimiento del minero}

Es el análogo a ``copiar 8~KiB + \texttt{mmu\_mapPage(TASK\_CODE, ...)} + \texttt{tlbflush}'' del TP x86:

\begin{codesnippet}
\begin{verbatim}
 void mmu_task_move(uint32_t slot,
                    uint32_t old_x, uint32_t old_y,
                    uint32_t new_x, uint32_t new_y) {
     TaskMmu& t = g_tasks[slot];
     if (!t.initialized) return;

     const uint64_t src_pa = board_cell_pa(old_x, old_y);
     const uint64_t dst_pa = board_cell_pa(new_x, new_y);

     memcpy((void*)dst_pa, (const void*)src_pa, kCellBytes);
     fence_i();                              // ifetch coherence

     map_user_code_pages(t, dst_pa);
     sfence_vma_addr(kUserBase);
     sfence_vma_addr(kUserBase + PAGE_SIZE);
 }
\end{verbatim}
\end{codesnippet}

Detalles técnicos:

\begin{itemize}
  \item \textbf{\texttt{fence.i}.} Después de escribir instrucciones en memoria, RISC-V no garantiza que el pipeline de
  instrucciones vea los bytes nuevos hasta una \texttt{fence.i} explícita~\cite{riscv-unpriv}. En x86 la coherencia
  I\$/D\$ es automática; en RISC-V no, y sin esta instrucción el minero podría ejecutar el código viejo después del
  movimiento (me pasó debuggeando: la tarea arrancaba, se ``movía'', y seguía ejecutando como si estuviera en la celda
  original).
  \item \textbf{\texttt{sfence.vma rs1=va, rs2=x0}.} Invalida sólo la entrada del TLB para esa VA, el equivalente del
  \texttt{invlpg} de x86. Es más barato que un flush global y alcanza porque el único cambio fue \texttt{TASK\_CODE}.
\end{itemize}

\subsection{ASIDs}

Sv39 permite incluir un \texttt{ASID} en \texttt{satp} para que el TLB no haga flush global al cambiar de address space.
Este port asigna:

\begin{itemize}
  \item \texttt{asid}$=$0 para idle.
  \item \texttt{asid}$=$slot+1 para mineros (1..20).
\end{itemize}

El \texttt{sfence.vma x0, x0} sólo se usa al inicializar una tarea; en runtime los cambios de address space se hacen con
\texttt{mmu\_switch()} que escribe \texttt{satp} (con su ASID) y hace un \texttt{sfence.vma} global (simplicidad $>$
eficiencia).
